\documentclass[12pt]{article}
\usepackage[utf8]{inputenc}
\usepackage[T1]{fontenc}
\usepackage{hyperref}
\title{Pancana Language Resources for K-12 Curriculum Development}
\author{Compiled by Claude Sonnet 5.5 (claude-sonnet-5-5) for Andr\'as Kornai}
\date{}
\begin{document}
\maketitle
\section{Introduction}
Pancana (ISO 639-3 pnp) is an Austronesian language of Buton Island, off the southeast coast of Sulawesi, Indonesia, in the Muna--Buton group. About 7,000 native speakers are cited. Web searches found only general reference pages and no dictionary, grammar, primer or school material. This dossier is therefore short.
\section{Primary Educational Resources}
None found. No primers, textbooks or school programmes for Pancana turned up in search results.
\section{Linguistic and Reference Materials}
\begin{itemize}
\item Wikipedia, Pancana language: general overview and classification. \url{https://en.wikipedia.org/wiki/Pancana_language}
\item Max Planck Institute numeral archive, Pancana (fieldwork by Jonathan McDowell, 2010): numerals only. \url{https://lingweb.eva.mpg.de/channumerals/Pancana.htm}
\end{itemize}
\section{Religious and Community Resources}
\begin{itemize}
\item IDRAP profile of the Pancana people (community background). \url{https://idrap.or.id/eng/detailIndigenous.php?ID=14}
\item Joshua Project, Pancana in Indonesia (demographic profile; scripture translation listed as started). \url{https://joshuaproject.net/people_groups/print/14291/ID}
\end{itemize}
\section{Digital and Audio Resources}
None found.
\section{Curriculum Development Context}
Documentation appears thin. Any curriculum would need new fieldwork, and the closely related Muna and Cia-Cia languages may offer comparative material. Academic grammars may exist that these searches did not surface.
\end{document}